% !TeX root = ../main.tex
\section{Conclusion}
\label{sec:conclusion}

This research presents a Multi-Condition Underwater Image Enhancement network based on Rectified Flow, referred to as RF-MCUIE, designed to mitigate severe underwater degradations such as color casts and structural blurring. The framework addresses training instabilities through an Adaptive Masking strategy, termed ADMK, which effectively filters low-quality supervision signals. Furthermore, physical priors involving background light and transmission maps are integrated as conditional guidance to ensure the restoration adheres to established optical mechanisms. By leveraging the Rectified Flow regime and a two-stage composite loss, the model establishes a straight-line transport path that facilitates high-fidelity generation within minimal inference steps. Comprehensive experiments on the LSUI, UIEB, and U45 benchmarks demonstrate that the proposed method achieves state-of-the-art performance across various objective metrics, including LPIPS, FID, and UCIQE. Ultimately, RF-MCUIE provides a reliable and efficient solution for bandwidth-limited underwater visual systems by balancing superior restoration quality with significantly reduced computational latency.
